\section{Pourquoi c'est un accélérateur}
\label{sec:accel}

\subsection{Définition formelle d'un accélérateur matériel}

Un \textbf{accélérateur matériel} est un circuit spécialisé qui exécute une fonction spécifique significativement plus rapidement et/ou plus efficacement (en énergie par opération) qu'un processeur généraliste exécutant la même fonction en logiciel. Cette définition implique trois critères que notre co-processeur SHA-256 satisfait pleinement : un gain de performance mesurable et substantiel, un mécanisme physique expliquant ce gain, et une intégration en tant que périphérique auxiliaire (co-processeur) plutôt qu'en remplacement du processeur principal.

Dans cette section, nous démontrons quantitativement et qualitativement que notre design répond à cette définition en analysant trois sources distinctes de gain de performance : le parallélisme spatial, les opérations à coût nul, et l'absence d'overhead d'infrastructure.

\subsection{Raison 1 : Parallélisme spatial}

La différence fondamentale entre un processeur et un accélérateur matériel réside dans la manière dont les opérations sont exécutées temporellement. Un processeur séquentiel, même pipeliné, exécute fondamentalement une instruction à la fois sur une ALU partagée. Pour calculer $T_1 = h + \Sigma_1(e) + Ch(e,f,g) + K_t + W_t$ en logiciel, le processeur doit exécuter séquentiellement : le calcul de $\Sigma_1(e)$ (3 rotations + 2 XOR = environ 7 instructions), le calcul de $Ch(e,f,g)$ (2 AND + 1 NOT + 1 XOR = 4 instructions), puis quatre additions successives. Au total, une seule ronde de SHA-256 nécessite environ 50 instructions sur un ARM Cortex-M4.

En contraste, notre accélérateur calcule \textbf{toutes ces opérations simultanément} dans des chemins combinatoires parallèles dédiés. Les fonctions $\Sigma_1(e)$, $Ch(e,f,g)$, $\Sigma_0(a)$, $Maj(a,b,c)$, et l'expansion $W_t$ sont calculées en parallèle dans des sous-circuits indépendants, et leurs résultats convergent vers les additionneurs en une fraction de la période d'horloge. Le résultat final est disponible après un seul cycle, représentant un gain de 50$\times$ par ronde.

La figure~\ref{fig:cpu_vs_fpga} illustre cette différence fondamentale avec une comparaison temporelle.

\begin{figure}[H]
\centering
\begin{tikzpicture}[
    box/.style={draw, rectangle, minimum height=0.4cm, font=\tiny, thick},
]
% CPU timeline
\node[font=\small\bfseries] at (-0.5, 2) {CPU (séquentiel)};
\foreach \i in {0,...,12} {
    \node[box, fill=red!20, minimum width=0.38cm] at (\i*0.42, 1.4) {};
}
\node[font=\tiny] at (2.5, 0.9) {$\sim$50 instructions par ronde};
\draw[{Stealth[length=2mm]}-{Stealth[length=2mm]}, thick] (0,0.7) -- (5.4,0.7);
\node[font=\tiny] at (2.7, 0.4) {$\sim$50 cycles};

% FPGA timeline
\node[font=\small\bfseries] at (-0.5, -0.5) {FPGA (parallèle)};
\node[box, fill=green!30, minimum width=2.5cm, minimum height=0.6cm] at (1.25, -1.1)
    {\footnotesize \textbf{1 seul cycle}};
\draw[{Stealth[length=2mm]}-{Stealth[length=2mm]}, thick] (0,-1.7) -- (2.5,-1.7);
\node[font=\tiny] at (1.25, -2) {1 cycle};

% Speedup arrow
\draw[-{Stealth[length=3mm]}, ultra thick, blue!60!black] (6.5, 1.4) -- (6.5, -1.1)
    node[midway, right, font=\small\bfseries, blue!60!black] {$\times$50};
\end{tikzpicture}
\caption{CPU : $\sim$50 instructions/ronde vs FPGA : 1 cycle/ronde}
\label{fig:cpu_vs_fpga}
\end{figure}

Ce parallélisme spatial est possible grâce à la nature reconfigurable du FPGA qui permet d'instancier des centaines de portes logiques fonctionnant simultanément, là où un processeur ne dispose que d'une ou deux ALU 32 bits partagées entre toutes les instructions.

\subsection{Raison 2 : Opérations à coût nul en matériel}

L'algorithme SHA-256 fait un usage intensif de rotations et de décalages de bits. Sur un processeur ARM Cortex-M4, chaque rotation nécessite 1 à 3 instructions (selon que le processeur dispose d'une instruction \texttt{ROR} native ou doit la simuler par shift + OR). Sur notre accélérateur FPGA, ces opérations sont \textbf{littéralement gratuites} : elles ne consomment aucune LUT, aucune bascule, et n'ajoutent aucun délai de propagation.

\begin{figure}[H]
\centering
\begin{tikzpicture}[
    box/.style={draw, rectangle, minimum width=3.5cm, minimum height=0.6cm,
                align=left, font=\footnotesize, thick},
]
\node[box, fill=red!15] (cpu) at (0,0) {CPU: ROTR = 2-3 instructions\\(shift + OR + mask)};
\node[box, fill=green!15] (fpga) at (0,-1.3) {FPGA: ROTR = 0 ressource\\(simple reroutage de fils)};
\node[font=\small\bfseries] at (-3, 0) {Rotation:};
\draw[-{Stealth[length=2mm]}, thick, blue] (cpu.south) -- (fpga.north)
    node[midway, right, font=\scriptsize, blue] {Gratuit !};
\end{tikzpicture}
\caption{Les rotations sont gratuites en matériel (reroutage de fils)}
\label{fig:rotr_free}
\end{figure}

L'explication physique est simple : en matériel, une rotation de $n$ positions correspond à un réarrangement des connexions entre les étages de registres. Les 32 fils de sortie du registre source sont simplement reconnectés aux 32 entrées du registre destination dans un ordre décalé. Ce réarrangement est résolu au moment du placement-routage par l'outil de synthèse et ne nécessite aucune porte logique active. Le « coût » se limite à la longueur des fils, qui est négligeable devant le délai d'une porte logique.

Quantitativement, SHA-256 effectue 48 rotations et 16 décalages par ronde (dans $\Sigma_0$, $\Sigma_1$, $\sigma_0$, $\sigma_1$), soit $64 \times 64 = 4096$ rotations/décalages par bloc. Sur un processeur, cela représente 4096 à 12288 instructions ; sur notre FPGA, cela représente exactement 0 cycle additionnel et 0 ressource logique.

De même, les opérations XOR sur 32 bits se traduisent en matériel par une seule couche de 32 portes XOR2 fonctionnant en parallèle, avec un délai de propagation inférieur à 0.5 ns (une seule LUT-level). La sélection de $K_t$ (la constante de ronde) est implémentée par un multiplexeur câblé indexé par le compteur --- essentiellement un tableau de lookup combinatoire --- plutôt que par une lecture mémoire coûteuse en latence.

\subsection{Raison 3 : Absence d'overhead d'infrastructure}

Un processeur généraliste, même lorsqu'il exécute un algorithme aussi simple que SHA-256, consacre une fraction significative de ses ressources et de son temps à l'infrastructure logicielle et matérielle :

\begin{itemize}
    \item \textbf{Fetch d'instructions} : chaque instruction doit être lue depuis le cache L1 ou la mémoire programme (1--3 cycles de latence par instruction en cas de miss) ;
    \item \textbf{Décodage} : l'instruction binaire doit être décodée pour déterminer l'opération à effectuer, les registres sources et destination (1 cycle de pipeline) ;
    \item \textbf{Gestion des registres} : avec seulement 13 registres généraux sur ARM Cortex-M4, les 8 variables de travail $a..h$ de SHA-256 plus les variables temporaires nécessitent des spills vers la pile (accès mémoire coûteux) ;
    \item \textbf{Appels de fonction et pile} : les sous-routines pour les fonctions $Ch$, $Maj$, etc. ajoutent des cycles de push/pop/branch ;
    \item \textbf{Interruptions et changements de contexte} : sur un système multitâche, le calcul peut être interrompu à tout moment.
\end{itemize}

Notre accélérateur élimine \textbf{intégralement} cet overhead. Il n'a pas d'instructions à fetcher (son comportement est câblé), pas de décodeur (les signaux de contrôle sont directement générés par la FSM), pas de limitation de registres (les 8 registres de travail plus les 16 registres du Scheduler sont tous instanciés physiquement), et il fonctionne de manière déterministe sans interruption. Chaque cycle d'horloge produit un travail utile directement lié au calcul du condensé : l'efficacité est de 100\%.

\subsection{Comparaison quantitative détaillée}

Le tableau~\ref{tab:perf} rassemble les mesures de performance pour différentes plateformes, permettant de situer notre accélérateur dans le paysage des solutions existantes.

\begin{table}[H]
\centering
\caption{Comparaison de performance : logiciel vs accélérateur matériel}
\label{tab:perf}
\begin{tabular}{lrrrr}
\toprule
\textbf{Plateforme} & \textbf{Cyc/bloc} & \textbf{Freq} & \textbf{Débit} & \textbf{Gain} \\
\midrule
ARM Cortex-M4 (logiciel optimisé) & 3\,500 & 168 MHz & 25 Mb/s & 1$\times$ \\
ARM Cortex-A53 (logiciel + NEON) & 1\,800 & 1.2 GHz & 341 Mb/s & 14$\times$ \\
\midrule
\textbf{Notre accel. (100 MHz)} & \textbf{85} & 100 MHz & \textbf{602 Mb/s} & \textbf{24$\times$} \\
Notre accel. (200 MHz) & 85 & 200 MHz & 1\,204 Mb/s & 48$\times$ \\
\midrule
McLoone \& McCanny~\cite{mcloone} (pipeliné) & 64 & 150 MHz & 1\,200 Mb/s & 48$\times$ \\
Dadda et al.~\cite{dadda} (ASIC) & 64 & 300 MHz & 2\,400 Mb/s & 96$\times$ \\
\bottomrule
\end{tabular}
\end{table}

Les calculs de débit sont effectués selon la formule :
\[
\text{Débit} = \frac{\text{Taille du bloc}}{\text{Cycles par bloc} \times \text{Période d'horloge}} = \frac{512 \text{ bits}}{N_{cyc} \times T_{clk}}
\]

Pour notre accélérateur à 100 MHz (période 10 ns) :
\[
\text{Débit} = \frac{512}{85 \times 10 \text{ ns}} = \frac{512}{850 \text{ ns}} = 602{,}4 \text{ Mbit/s}
\]

Le débit asymptotique (pour un message long de $N$ blocs, où l'overhead de WRITEBACK est amorti) est plus élevé car le coût marginal par bloc est de 67 cycles plutôt que 85 :
\[
\text{Débit}_{asymp} = \frac{512}{67 \times 10 \text{ ns}} = 764 \text{ Mbit/s}
\]

\subsection{Analyse du gain par rapport au Cortex-M4}

Le gain de 24$\times$ par rapport au Cortex-M4 se décompose comme suit :

\begin{itemize}
    \item \textbf{Réduction du nombre de cycles} : $3500 / 85 = 41\times$ (le facteur principal) ;
    \item \textbf{Pénalité de fréquence} : notre accélérateur à 100 MHz est cadencé $168/100 = 1.68\times$ plus lentement que le Cortex-M4 ;
    \item \textbf{Gain net} : $41 / 1.68 \approx 24\times$.
\end{itemize}

Si l'accélérateur est synthétisé à 200 MHz (atteignable sur Artix-7 avec des optimisations du chemin critique), le gain monte à $41 \times 200/168 \approx 48\times$. Ce gain est obtenu avec une surface de l'ordre de 1\,500--2\,000 LUT, soit moins de 4\% des ressources d'un Artix-7 modeste (XC7A35T).

\subsection{Comparaison avec les publications académiques}

Notre architecture itérative à une ronde par cycle se compare favorablement avec les designs publiés dans la littérature pour des architectures similaires :

\begin{itemize}
    \item \textbf{McLoone et McCanny (IEEE FPL 2002)}~\cite{mcloone} présentent une architecture pipelinée à 64 étages pour SHA-384/512. Leur débit est plus élevé (1.2 Gb/s) mais au prix d'une surface 64 fois supérieure, inadaptée à un co-processeur embarqué.
    \item \textbf{Dadda et al. (GLSVLSI 2004)}~\cite{dadda} proposent une implémentation ASIC haute performance utilisant des additionneurs carry-save pour réduire le chemin critique. Leur fréquence de 300 MHz sur ASIC est supérieure à notre cible FPGA, mais leur approche n'est pas reconfigurable.
    \item \textbf{Dowenbergh et al. (ICACT 2019)}~\cite{icact2019} présentent un cluster FPGA-CPU pour SHA-256, démontrant que l'approche co-processeur est pertinente même à grande échelle.
\end{itemize}

Notre design se situe dans la catégorie « co-processeur compact et efficace » : il offre un excellent ratio débit/surface pour les applications embarquées où la surface disponible est limitée.

\subsection{Facteurs limitant la Fmax}

La fréquence maximale de notre design est limitée par le chemin critique qui traverse le module Round. Ce chemin comprend :

\begin{enumerate}
    \item La lecture du registre $h$ (délai $t_{cko}$ du flip-flop : $\sim$0.5 ns) ;
    \item Le calcul de $\Sigma_1(e)$ (1 niveau de XOR : $\sim$0.5 ns) ;
    \item Le calcul de $Ch(e,f,g)$ (1 niveau de LUT : $\sim$0.5 ns) ;
    \item Quatre additions de 32 bits en cascade ($4 \times 2.5$ ns $= 10$ ns avec carry chains) ;
    \item Le temps de setup du flip-flop destination ($t_{su}$ : $\sim$0.2 ns).
\end{enumerate}

Le total théorique est d'environ 12 ns (83 MHz). En pratique, l'outil de synthèse optimise les additions en réarrangeant l'arbre d'opérandes (carry-save partial products) et en exploitant les primitives CARRY4 dédiées du FPGA, permettant d'atteindre 5--7 ns de chemin critique (143--200 MHz). Une optimisation classique pour repousser davantage cette limite consiste à pré-calculer $K_t + W_t$ un cycle à l'avance dans un registre dédié, réduisant la chaîne d'addition de 5 termes à 4 termes et gagnant environ 20--25\% de Fmax.

\subsection{Compromis surface-débit et analyse d'efficacité}

L'efficacité d'un accélérateur peut être mesurée par le ratio débit/surface, exprimé en Mbit/s par LUT. Pour notre design :

\[
\text{Efficacité} = \frac{602 \text{ Mbit/s}}{{\sim}1800 \text{ LUT}} \approx 0.33 \text{ Mbit/s/LUT}
\]

Une architecture pipelinée à 64 étages offrirait un débit de $512 \times 200\text{ MHz} / 1 = 102\,400$ Mbit/s théoriques, mais avec $64 \times 1800 = 115\,200$ LUT, soit une efficacité de $102400/115200 \approx 0.89$ Mbit/s/LUT. Le pipeline est donc plus efficace par LUT (facteur 2.7$\times$) mais exige 64 fois plus de surface --- une ressource qui n'est tout simplement pas disponible sur un FPGA embarqué si l'on souhaite également implémenter un processeur et des périphériques.

\subsection{Synthèse visuelle du gain d'accélération}

La figure~\ref{fig:accel_gain} résume visuellement le bénéfice de l'approche co-processeur pour le système complet. Non seulement le calcul est 24 fois plus rapide, mais le processeur hôte est libéré pendant toute la durée du hachage, lui permettant d'exécuter d'autres tâches (communication réseau, gestion de protocoles, interface utilisateur).

\begin{figure}[H]
\centering
\begin{tikzpicture}[
    block/.style={draw, rectangle, rounded corners, minimum width=4cm,
                  minimum height=1.5cm, align=center, font=\footnotesize, thick},
]
% Left: CPU alone
\node[block, fill=red!10] (cpu) at (0,0) {
    \textbf{CPU seul}\\[3pt]
    3\,500 cycles/bloc\\
    CPU bloqué pendant le calcul\\
    Débit : 25 Mbit/s
};

% Right: CPU + accelerator
\node[block, fill=green!10] (accel) at (7,0) {
    \textbf{CPU + Accélérateur}\\[3pt]
    85 cycles/bloc\\
    CPU libre pour autres tâches\\
    Débit : 602 Mbit/s
};

% Arrow
\draw[-{Stealth[length=4mm]}, ultra thick, blue!60!black] (cpu.east) -- (accel.west)
    node[midway, above, font=\large\bfseries, blue!60!black] {$\times$24};
\end{tikzpicture}
\caption{Gain d'accélération : CPU seul vs CPU + co-processeur SHA-256}
\label{fig:accel_gain}
\end{figure}

\subsection{Conclusion : un véritable accélérateur}

En résumé, notre co-processeur SHA-256 satisfait pleinement la définition d'un accélérateur matériel :

\begin{enumerate}
    \item \textbf{Gain de performance substantiel} : 24--48$\times$ selon la fréquence, mesuré et justifié par des calculs vérifiables ;
    \item \textbf{Mécanismes physiques identifiés} : parallélisme spatial (50$\times$ par ronde), opérations à coût nul (rotations/XOR), absence d'overhead d'instruction (100\% d'efficacité par cycle) ;
    \item \textbf{Architecture de co-processeur} : le design fonctionne en tant que périphérique esclave du processeur hôte, communiquant via une interface mémoire simple, et libérant le CPU pendant le calcul.
\end{enumerate}

Le débit de 602 Mbit/s à 100 MHz est suffisant pour saturer une interface Gigabit Ethernet en temps réel (le SHA-256 des paquets TLS peut être calculé au fil de l'eau), ce qui démontre la pertinence pratique de l'accélération matérielle pour les applications réseau sécurisées.

